% Part: first-order-logic
% Chapter: proof-systems
% Section: axiomatic-deduction

\documentclass[../../../include/open-logic-section]{subfiles}

\begin{document}

\iftag{FOL}
      {\olfileid{fol}{prf}{axd}}
      {\olfileid{pl}{prf}{axd}}

\olsection{\usetoken{P}{derivation} Aksiomatik}

!!{derivation}s aksiomatik iku sistem !!{derivation} logis kang paling
tuwa lan paling prasaja. !!{derivation}s ing sistem iki mung arupa
urutan !!{sentence}s. Urutan !!{sentence}s dianggep minangka
!!{derivation} kang bener yen saben !!{sentence}~$!A$ ing jerone
nyukupi salah siji syarat iki:
\begin{enumerate}
\item $!A$ iku aksioma, utawa
\item $!A$ iku !!a{element} saka himpunan~$\Gamma$ saka !!{sentence}s
  kang diwenehake, utawa
\item $!A$ dibenerake dening aturan inferensi.
\end{enumerate}
Supaya dadi aksioma, $!A$ kudu duwe wujud salah siji saka sawetara
skema !!{sentence} kang tetep. Ana akeh himpunan skema aksioma kang
menehi sistem !!{derivation} kang nyukupi (sah lan jangkep) kanggo
logika tataran kapisan. Sawetara ditata miturut konektif kang diatur,
upamane skema
\[
!A \lif (!B \lif !A) \qquad !B \lif (!B \lor !C) \qquad (!B \land !C) \lif !B
\]
iku aksioma kang umum kanggo ngatur $\lif$, $\lor$, lan~$\land$.
Sawetara sistem aksioma ngupaya cacah aksioma kang paling sithik.
Gumantung marang konektif kang dianggep primitif, kita malah bisa
nemokake sistem aksioma kang mung dumadi saka siji aksioma.

Aturan inferensi iku pratelan kondisional kang menehi syarat cukup
supaya !!a{sentence} ing !!a{derivation} bisa dibenerake. Modus ponens
iku salah siji aturan kaya mangkono kang umum banget: aturan iki kandha
yen $!A$ lan $!A \lif !B$ wis dibenerake, mula $!B$ uga dibenerake. Iki
ateges larik ing !!a{derivation} kang ngemot !!{sentence}~$!B$
dibenerake, angger $!A$ lan $!A \lif !B$ kalorone (kanggo sawijining
!!{sentence}~$!A$) muncul ing !!{derivation} sadurunge~$!B$.

Relasi $\Proves$ adhedhasar !!{derivation}s aksiomatik ditemtokake kaya
mangkene: $\Gamma \Proves !A$ yen lan mung yen ana !!a{derivation}
kang nduweni !!{sentence}~$!A$ minangka formula pungkasan (lan
$\Gamma$ dianggep minangka himpunan !!{sentence}s ing !!{derivation}
iku kang dibenerake dening~(2) ing ndhuwur). $!A$ iku teorema yen~$!A$
duwe !!a{derivation} nalika~$\Gamma$ kosong, yaiku saben !!{sentence}
ing !!{derivation} dibenerake dening (1) utawa~(3). Tuladhane, iki
!!a{derivation} kang nuduhake yen $\Proves
!A  \lif (!B \lif (!B \lor !A))$:
\begin{derivation}
  1. & $!B \lif (!B \lor !A)$ \\
  2. & $(!B \lif (!B \lor !A)) \lif (!A  \lif (!B \lif (!B \lor !A)))$\\
  3. & $!A  \lif (!B \lif (!B \lor !A))$
\end{derivation}
!!{sentence} ing larik~1 nduweni wujud aksioma $!A \lif (!A
\lor !B)$ (kanthi peran $!A$ lan $!B$ diijolake). Ukara ing larik~2
nduweni wujud aksioma $!A \lif (!B \lif !A)$. Mula, rong larik kasebut
dibenerake. Larik~3 dibenerake dening modus ponens: yen dicekakake
minangka $!D$, larik~2 duwe wujud $!C \lif !D$, kanthi $!C$ yaiku
$!B \lif (!B \lor !A)$, yaiku larik~1.

Himpunan $\Gamma$ ora konsisten yen $\Gamma \Proves \lfalse$. Sistem
aksioma kang jangkep uga bakal mbuktekake yen $\lfalse \lif !A$ kanggo
saben~$!A$, mula yen $\Gamma$ ora konsisten, $\Gamma \Proves !A$
kanggo saben~$!A$.

Sistem !!{derivation}s aksiomatik kanggo logika sepisanan diwenehake
dening Gottlob Frege ing \emph{Begriffsschrift} taun 1879, kang amarga
iku kerep dianggep minangka karya kapisan logika modern. Sistem iki
dadi sampurna ing \emph{Principia Mathematica} anggitane Alfred North
Whitehead lan Bertrand Russell lan liwat pakaryane David Hilbert lan
para muride ing taun 1920-an. Mula sistem kasebut kerep diarani
``sistem Frege'' utawa ``sistem Hilbert''. Sistem iki luwes banget
amarga nemokake sistem aksiomatik kanggo sawijining logika kerep
gampang. Amarga !!{derivation}s duwe struktur kang prasaja banget lan
mung siji utawa rong aturan inferensi, mbuktekake bab-bab
\emph{ngenani} derivasi kasebut uga cukup gampang. Nanging, sistem iki
angel banget digunakake ing praktik, yaiku angel nemokake lan nulis
bukti.

\end{document}
